\section{Shared insights across contributions}
\label{sec:disc_shared_insights}

Although the three main contributions of this thesis address challenges at different stages of the life cycle of a docked \ac{bss}, they share common underlying mechanisms. This section synthesizes these shared insights and reflects on their relevance for research and operational practice in bike-sharing.

A central theme across all three contributions is the importance of high-resolution and high-quality data for meaningful analysis and prediction in \acp{bss}. Across planning and operations, the results show that spatiotemporally detailed, contextual and operational data can reveal pronounced heterogeneities within the system. Different parts of the city, stations, and bicycles exhibit distinct patterns of accessibility, usage, energy depletion, and component wear that would be largely obscured under aggregated representations. Moreover, beyond supporting descriptive analyses, higher-resolution data can also lead to improved predictions by reducing the need for simplifying assumptions.

Another important topic concerns the role of network effects in shaping system behavior. The analyses consistently indicate that outcomes depend not only on the local attributes of stations or trips, but also on their position within the broader cycling network and their relationships to other system elements. Inter-station proximity and system accessibility influence feasible station layouts; mobility flows are constrained by the structure of the cycling network rather than by simple distance alone; and concentrated usage of specific network connections contributes to uneven battery usage and component degradation. These observations reinforce the need to treat \acp{bss} as networked systems, where spatial structure and connectivity determine how local characteristics translate into system-wide patterns. Furthermore, capturing these effects accurately requires not only modeling network structure based on distance, but also taking into account topographical variations and the directionality of travel.

Within this networked context, structural design choices—particularly station placement—are likely to shape operational dynamics. While this thesis does not explicitly quantify causal links between planning decisions and operational outcomes, the results suggest plausible pathways through which station layouts influence mobility patterns, which in turn affect battery dynamics and maintenance needs. This perspective underscores the limitations of treating planning, rebalancing, and maintenance as independent problems and highlights the value of viewing these components as interconnected parts of a single system. At the same time, user behavior remains adaptive and responsive to system changes, which helps explain why observed patterns may evolve over time.

In parallel, ongoing technological developments are expanding both the scope and the practical relevance of data-driven approaches in \acp{bss}. New generations of shared bicycles increasingly incorporate \ac{gps} units and a growing range of onboard sensors, enabling continuous monitoring of location, usage, and mechanical conditions. These include signals such as tire pressure and vibration that can provide additional information on bicycle state and deterioration. Meanwhile, contextual information is also becoming more accessible through the expansion of open urban data and more widespread data-sharing practices among public authorities. Together, these trends are lowering barriers to adoption and increasing operator interest in computational methods to improve efficiency, reduce costs, and optimize resource allocation. As a result, although some analyses developed in this thesis—particularly those related to battery dynamics and component wear—remain relatively uncommon in the current literature, growing data availability is likely to make such approaches progressively more feasible in real-world deployments.

Beyond methodological contributions, the findings have implications for \ac{bss} users and operators, and for society more broadly. While the analyses are grounded in Barcelona and its \ac{bss}, the proposed methods are largely transferable to other urban contexts, provided that comparable spatial, operational, and contextual data are available. For users, improved station placement and more reliable vehicle availability can translate into more consistent service quality. Planning approaches that explicitly incorporate accessibility and equity considerations can also support more balanced spatial coverage, helping to extend the benefits of bike-sharing to a wider range of neighborhoods and user groups. For operators, data-driven planning and predictive tools create opportunities to allocate resources more efficiently and reduce operational costs. At a societal level, more efficient, equitable, and reliable \acp{bss} can strengthen their role as a sustainable urban mobility option by supporting modal shift away from private motorized transport and contributing to lower emissions and more livable cities. In this sense, the approaches presented in this thesis align operational efficiency with broader environmental and social objectives.

Finally, these shared insights point to several directions for future research. A first and practical next step is to move beyond offline validation and implement the proposed methods in production environments. Deploying these approaches within operator data pipelines and decision workflows would enable iterative refinement based on real-world feedback. This includes feedback from planners on station-location proposals and from operators on battery status and component failure forecasts. To adapt to evolving usage patterns and system changes, these methods could also be continuously updated using streaming data. Building on this initial deployment experience, further improvements could leverage richer sensor data (e.g., \ac{gps} trajectories and onboard condition monitoring) to refine the proposed models. In addition, future deployments could evaluate these methods not only in terms of predictive accuracy but also in terms of operational value, using metrics such as service-level reliability, avoided empty stations, reduced downtime, and cost and emissions impacts.

At a broader level, a natural next step is to develop integrated decision-support frameworks that jointly consider station planning, battery-aware operations, and predictive maintenance under shared resource and service-level constraints, rather than optimizing these components in isolation. Empirically, this motivates works that quantify causal links between structural design choices and downstream operational outcomes, for instance by leveraging natural experiments around network expansions or station relocations. Finally, applying and benchmarking these approaches across multiple cities would help assess transferability and establish clearer evidence on how data quality, urban form, and topography shape model performance.
